\documentclass[11pt,a4paper]{article}
\usepackage[utf8]{inputenc}
\usepackage[T1]{fontenc}
\usepackage{amsmath, amssymb}
\usepackage{booktabs}
\usepackage[margin=2.5cm]{geometry}
\usepackage{hyperref}
\hypersetup{pdftitle={Energy Decay in a Damped Harmonic Oscillator}, pdfauthor={Lena Ortiz}}

\title{Energy Decay in a Damped Harmonic Oscillator}
\author{Lena Ortiz\\ \small Department of Physics, Northfield College}
\date{October 2026}

\begin{document}
\maketitle

\begin{abstract}
We study the free motion of a mass on a linear spring subject to viscous damping. We derive the closed-form solution in the underdamped regime, show that the mechanical energy decays monotonically, and compare the analytic solution with two simple numerical schemes.
\end{abstract}

\section{Introduction}\label{sec:intro}

The damped harmonic oscillator is the simplest model of a system that loses energy to its surroundings~\cite{ref1}. Its equation of motion is
\begin{equation}
  m\ddot{x}+c\dot{x}+kx=0
\end{equation}
where $m>0$ is the mass, $c\geq0$ the damping coefficient and $k>0$ the spring constant. Dividing by $m$ gives the standard form
\begin{equation}
  \ddot{x}+2\gamma\dot{x}+\omega_{0}^{2x}=0
\end{equation}

\section{Solution}\label{sec:method}

Substituting $x=e^{\lambda}t$ into (2) gives the characteristic equation $\lambda^{2}+2\gamma\lambda+\omega_{0}^{2}=0$, whose roots are
\begin{equation}
  \lambda_{\pm}=-\gamma\pm\sqrt{\gamma^{2}-\omega_{0}^{2}}
\end{equation}
The behaviour depends on the sign of the discriminant:
\begin{align}
  regime &=\begin{cases}
underdamped, & \gamma<\omega_{0}, \\
critically & damped,
\end{cases}\gamma=\omega_{0}, \\
  overdamped,
\end{align}

\section{Numerical check}\label{sec:results}

We integrated (2) with $\omega_{0}=2$, $\gamma=0.2$ and $x_{0}=1$ on $t\in[0,10]$, using the explicit Euler method and the classical fourth-order Runge--Kutta method (RK4)~\cite{ref2}. For each step size $h$ we:
\begin{enumerate}
  \item integrated from t = 0 to t = 10 with 10/h steps;
  \item compared every step with the exact solution (4);
  \item recorded the largest absolute difference.
\end{enumerate}

\begin{table}[htbp]
  \centering
  \caption{Table 1: Maximum absolute error on [0, 10].}
  \label{tab:results}
  \begin{tabular}{lcc}
    \toprule
    \textbf{Quantity} & \textbf{Value} & \textbf{Uncertainty} \\
    \midrule
    a & 1.00 & 0.02 \\
    b & 2.50 & 0.05 \\
    \bottomrule
  \end{tabular}
\end{table}

\begin{thebibliography}{9}
\bibitem{ref1} H. Goldstein, C. Poole and J. Safko, \emph{Classical Mechanics}, 3rd ed., Addison-Wesley, 2002.
\bibitem{ref2} E. Hairer, S. P. Nørsett and G. Wanner, \emph{Solving Ordinary Differential Equations I}, 2nd ed., Springer, 1993.
\end{thebibliography}

\end{document}
